% Unidad U016. Registro firmado, lectura terminal y reconstruccion de K.
% Redaccion autonoma desde la autoridad material 1063 y su sucesor V2.

\chapter{Registro dodecafásico firmado y reconstrucción integral del sello}
\label{chap:registro-dodecafasico-hadamard-k}

La construcción del sello entero no comienza con \(\alpha\), con su inversa
ni con una tabla metrológica. Su dominio es el estado terminal enriquecido
del Topological Phase Kernel obtenido después de la cadena ya construida
\begin{equation}
 \boxed{
 \mathrm{APP}\longrightarrow\mathrm{TRIT}\longrightarrow\mathrm{TPK}
 \longrightarrow\Omega_{\rm seed}
 \xrightarrow{T_{\min}}\mathcal U_6^{\rm cat}
 \xrightarrow{\rho_3}W_{243}
 \longrightarrow w_{30}\longrightarrow R_{36}
 \longrightarrow G_9\longrightarrow x_{\rm term}.}
 \label{eq:cadena-upstream-registro-dodecafasico}
\end{equation}
Las flechas anteriores fueron definidas con sus respectivos dominios:
pares de semillas, firmas, emisiones, palabras ternarias, fronteras,
historias compatibles y holonomía nonádica. Esta unidad no vuelve a escoger
ninguno de esos datos. Construye las cartas enteras del mismo estado
terminal y demuestra su equivalencia.

\section{Ventanas de la traza orientada de longitud ciento ocho}

Sea
\[
 \Theta_{108}=\mathbb Z/108\mathbb Z
\]
el calendario observable. Sus doce ventanas nonádicas son
\[
 E_m=\{9m+1,\ldots,9m+9\},
 \qquad 0\leq m<12.
\]
La partición es disjunta y exhaustiva:
\[
 \{1,\ldots,108\}=\bigsqcup_{m=0}^{11}E_m.
\]
El conjunto de códigos de dirección que publican un evento se fija como
\[
 D_{\rm evt}=\{2,4,6,8\}.
\]
La orientación de las doce ventanas es
\begin{equation}
 \Sigma_{12}=(+,+,+,-,-,-,+,+,+,-,-,-).
 \label{eq:signatura-doce-ventanas}
\end{equation}
Si \(d_t\) es la dirección codificada en el paso \(t\), el evento firmado
de la ventana \(m\) es
\begin{equation}
 e_m=\Sigma_{12}(m)\,
 \#\{t\in E_m:d_t\in D_{\rm evt}\}.
 \label{eq:evento-firmado-ventana}
\end{equation}
La fórmula distingue el número no negativo de eventos y el signo de la
orientación. No reemplaza una dirección cardinal por un signo: las dos
coordenadas siguen presentes en el estado TPK.

\begin{definition}[Registro enriquecido dodecafásico]
\label{def:registro-enriquecido-dodecafasico}
El registro transportado por la traza es
\[
 \mathcal R_{12}(x)=
 \bigl((E_m,\tau_m,\sigma_m,\kappa_m,\Lambda_m)\bigr)_{m=0}^{11},
\]
donde \(\tau_m\) es la orientación TRIT, \(\sigma_m\) la firma de
frontera, \(\kappa_m\) el acarreo, y \(\Lambda_m\) el segmento de ledger
que conserva la ruta. Existe una cadena de proyecciones
\[
 \mathcal X_{\rm TPK}^{\rm enr}
 \longrightarrow\mathcal R_{12}
 \longrightarrow\Theta_{108},
\]
pero ninguna de las dos flechas es una identificación.
\end{definition}

En particular, \(\Theta_{108}\) sólo conserva la posición del calendario.
El registro \(\mathcal R_{12}\) conserva además orientación, firma,
acarreo y procedencia. El estado completo conserva todavía cilindro,
frontera, cociente, terminal y memoria vertical.

\section{Carta integral de la oposición dodecafásica}

Para \(0\leq i<6\), las posiciones opuestas determinan
\begin{equation}
 p_i=e_i+e_{i+6},
 \qquad
 a_i=e_i-e_{i+6}.
 \label{eq:pares-opuestos-registro}
\end{equation}
La carga central y los cinco márgenes relativos son
\begin{equation}
 s_C=\sum_{i=0}^{5}p_i,
 \qquad
 M_i=p_i-p_5\quad(0\leq i<5).
 \label{eq:carga-margenes-registro}
\end{equation}

\begin{definition}[Carta integral \(1+5+6\)]
\label{def:carta-integral-156}
Se define
\begin{equation}
 \mathcal L_{12}(e_0,\ldots,e_{11})
 =\bigl(s_C,M_0,\ldots,M_4,a_0,\ldots,a_5\bigr).
 \label{eq:carta-integral-156}
\end{equation}
La partición de coordenadas es
\[
 1+5+6=12.
\]
\end{definition}

\begin{theorem}[Inversión exacta de la carta integral]
\label{thm:inversion-carta-integral-156}
Una tupla
\[
 \ell=(s_C,M_0,\ldots,M_4,a_0,\ldots,a_5)\in\mathbb Z^{12}
\]
pertenece a la imagen de \(\mathcal L_{12}\) si y sólo si
\begin{align}
 s_C-\sum_{i=0}^{4}M_i&\equiv0\pmod6,
 \label{eq:congruencia-seis-carta}\\
 p_i&\equiv a_i\pmod2
 \quad(0\leq i<6),
 \label{eq:congruencia-paridad-carta}
\end{align}
donde
\begin{align}
 p_5&=\frac{s_C-\sum_{i=0}^{4}M_i}{6},
 \label{eq:inversa-p5-carta}\\
 p_i&=p_5+M_i\quad(0\leq i<5).
 \label{eq:inversa-pi-carta}
\end{align}
En ese caso, la única preimagen entera es
\begin{equation}
 e_i=\frac{p_i+a_i}{2},
 \qquad
 e_{i+6}=\frac{p_i-a_i}{2}.
 \label{eq:inversa-eventos-carta}
\end{equation}
\end{theorem}

\begin{proof}
Sumando \(M_i=p_i-p_5\) para \(0\leq i<5\), obtenemos
\[
 \sum_{i=0}^{4}M_i
 =\sum_{i=0}^{4}p_i-5p_5
 =s_C-6p_5,
\]
lo que fuerza \eqref{eq:inversa-p5-carta} y la congruencia módulo seis.
Una vez reconstruidos los seis \(p_i\), el sistema
\[
 p_i=e_i+e_{i+6},
 \qquad a_i=e_i-e_{i+6}
\]
tiene la única solución \eqref{eq:inversa-eventos-carta}. Esa solución es
entera exactamente cuando \(p_i\) y \(a_i\) tienen la misma paridad.
Las dos familias de congruencias son, por tanto, necesarias y suficientes.
\end{proof}

Este teorema es anterior al cierre numérico de \(\alpha\). Explica cómo la
traza orientada produce una carta integral reversible, y por qué no basta
con conservar el calendario reducido.

\section{Selección del estado terminal común}

Las tres secciones compatibles llegan a profundidad terminal con catálogos
de cardinalidades
\[
 |\mathcal T_\pi|=196,
 \qquad
 |\mathcal T_e|=197,
 \qquad
 |\mathcal T_\varphi|=196.
\]
El producto contiene
\begin{equation}
 196\cdot197\cdot196=7\,567\,952
 \label{eq:cardinal-producto-terminal}
\end{equation}
ternas ordenadas.

El margen de columnas del lector terminal es
\[
 q_\partial=(1,2,2,1,2,0)\in\mathbb F_3^6,
\]
y la orientación de filas del canal conjunto es
\[
 \sigma=(1,1,-1)=(1,1,2)\in\mathbb F_3^3.
\]
El primer dato actúa sobre columnas; el segundo sobre filas. No se deduce
uno del otro.

Las firmas locales de Gram, norma, peso y distancia reducen
\eqref{eq:cardinal-producto-terminal} a \(331\) ternas. El margen
\(q_\partial\) deja dos índices:
\[
 (120,197,157),
 \qquad
 (129,197,148).
\]
Sus matrices son
\[
 B_0=
 \begin{pmatrix}
 0&2&0&2&2&0\\
 0&0&1&2&2&0\\
 1&0&1&0&1&0
 \end{pmatrix},
 \qquad
 B_1=
 \begin{pmatrix}
 0&2&1&0&2&0\\
 0&0&1&2&2&0\\
 1&0&0&2&1&0
 \end{pmatrix}.
\]
Las sumas de filas son
\[
 (0,2,0),
 \qquad
 (2,2,1).
\]
Como
\[
 \langle\sigma\rangle
 =\{(0,0,0),(1,1,2),(2,2,1)\},
\]
sólo \(B_1\) satisface la condición axial.

\begin{theorem}[Unicidad de la publicación terminal visible]
\label{thm:unicidad-publicacion-terminal-visible}
La selección conjunta determina
\begin{equation}
 B_{\rm term}=
 \begin{pmatrix}
 021020\\001220\\100210
 \end{pmatrix},
 \label{eq:matriz-terminal-visible}
\end{equation}
correspondiente a la terna \((129,197,148)\). Ninguna evaluación decimal de
\(\alpha\) interviene en el censo ni en las dos condiciones de selección.
\end{theorem}

\begin{proof}
El censo y el margen dejan exactamente \(B_0,B_1\). La carga de \(B_0\)
no pertenece a \(\langle\sigma\rangle\), mientras la de \(B_1\) es
\(2\sigma\). La segunda candidata es, por tanto, la única que satisface
simultáneamente ambos requisitos tipados.
\end{proof}

\section{Dos publicaciones del mismo estado, no una conversión entre ellas}

Sea
\[
 x_{\rm term}\in
 \mathcal X_{\rm TPK}^{\rm term,enr}
 \subseteq\mathcal X_{\rm TPK}^{[108]}.
\]
La lectura visible es
\[
 \pi_{\rm term}(x_{\rm term})=B_{\rm term}.
\]
La carta firmada conserva otra coordenada del mismo estado.

\begin{definition}[Lector firmado producido por el registro terminal]
\label{def:subespacio-firmado-terminal}
El estado de entrada del lector es exclusivamente el estado terminal
enriquecido ya producido por la tubería causal
\(\mathcal U_t=\mathsf{Upd}_t\circ\mathsf{Tra}_t\circ\mathsf{Sel}_t\):
\begin{equation}
 x_{\rm term}
 =\bigl(\mathcal U_{108}\circ\cdots\circ\mathcal U_1\bigr)
   (x_{\rm seed})
 \in\mathcal X_{\rm TPK}^{\rm term,enr}.
 \label{eq:subespacio-firmado-terminal}
\end{equation}
La composición del registro orientado con sus lectores de canal \(90/120\)
extrae del libro \(\lambda_{108}\), sin añadir coordenadas externas,
\begin{equation}
 \mathcal E_{108}^{90,120}\!\left(\mathcal R_{12}(x_{\rm term})\right)
 =\bigl(b^{90},b^{120},Q_{\rm TPK}\bigr),
 \label{eq:canales-ledger-previos-u}
\end{equation}
donde
\begin{align}
 b^{90}={}&(-495,-116,-684,108,601,-90,
             -173,-88,313,560,-397,461),\notag\\
 b^{120}={}&(-425,-281,-481,671,-255,30,
              475,-543,680,251,6,-128),\notag\\
 Q_{\rm TPK}={}&6263.\notag
\end{align}
Estas veinticinco coordenadas son salidas del registro de la
\cref{def:registro-enriquecido-dodecafasico}, no datos añadidos al dominio.

Sobre una salida compatible \((b^{90},b^{120},Q)\), definimos el lector
armónico integral \(\Pi_H\) de forma material. Para \(r=1,2,3\), sean
\begin{equation}
 B_r=\sum_{j=0}^{3}b^{120}_{r+3j},
 \quad
 A_1=\frac{Q+2B_1+B_2}{3},
 \quad A_2=A_1-B_1,
 \quad A_3=A_2-B_2,
 \label{eq:lector-armonico-a-desde-ledger}
\end{equation}
y \(d_{r,j}=b^{90}_{r+3j}\). El bloque orbital \(r\) es
\begin{equation}
 \Pi_H^{(r)}(b^{90},b^{120},Q)
 =\bigl(A_r,d_{r,1}-d_{r,3},
              d_{r,0}+d_{r,2},d_{r,0}-d_{r,2}\bigr).
 \label{eq:lector-armonico-bloque-desde-ledger}
\end{equation}
Intercalando por columnas los tres bloques se obtiene \(\Pi_H\), y la
lectura firmada queda construida como la composición
\begin{equation}
 \boxed{
 R_{\rm sgn}:=
 \Pi_H\circ\mathcal E_{108}^{90,120}\circ\mathcal R_{12}:
 \mathcal X_{\rm TPK}^{\rm term,enr}
 \longrightarrow\mathcal U_{12}^{\rm sgn}:=
 \operatorname{im}R_{\rm sgn}\subseteq\mathbb Z^{12}.}
 \label{eq:lector-firmado-terminal}
\end{equation}
Ni \(U\) ni \(K\) pertenecen al dominio: \(U\) es la salida de
\eqref{eq:lector-firmado-terminal}, y \(K\) se reconstruye después mediante
la inversa integral de Hadamard.
\end{definition}

La escritura conjunta correcta es
\begin{equation}
 x_{\rm term}
 \xrightarrow{(\pi_{\rm term},R_{\rm sgn})}
 (B_{\rm term},U_{12}^{\rm sgn}).
 \label{eq:doble-publicacion-terminal}
\end{equation}
No se interpone una flecha
\(B_{\rm term}\to U_{12}^{\rm sgn}\): la matriz visible ha olvidado
precisamente las coordenadas de orientación, oposición, cociente, acarreo y
ledger utilizadas por la carta firmada. La fuente de ambas publicaciones es
el estado enriquecido producido por \eqref{eq:cadena-upstream-registro-dodecafasico}.

La evaluación del registro firmado da
\begin{equation}
 \boxed{
 U=(2378,1406,2479,-452,998,-551,
 -668,-204,-371,-322,-28,-997).}
 \label{eq:vector-u-firmado-autonomo}
\end{equation}

\section{Naturalidad respecto de la profundidad}

Para \(N'\geq N\geq108\), el truncamiento
\begin{equation}
 \tau_{N',N}:
 \mathcal X_{\rm TPK}^{[N'],\rm enr}
 \longrightarrow\mathcal X_{\rm TPK}^{[N],\rm enr}
 \label{eq:truncamiento-preserva-ku-autonomo}
\end{equation}
actúa sobre la historia enriquecida y su libro, no sobre una pareja
\((K,U)\) añadida al estado.

\begin{theorem}[Cuadrado de naturalidad de la lectura firmada]
\label{thm:naturalidad-lector-firmado-autonomo}
Se cumple
\begin{equation}
 \boxed{
 R_{{\rm sgn},N}\circ\tau_{N',N}
 =R_{{\rm sgn},N'}.}
 \label{eq:cuadrado-naturalidad-lector-firmado}
\end{equation}
\end{theorem}

\begin{proof}
Cada actualizador \(\mathcal U_t\) añade la incidencia siguiente al libro
sin reescribir las ciento ocho ventanas ya generadas. Por ello
\(\mathcal R_{12}\), \(\mathcal E_{108}^{90,120}\) y \(\Pi_H\) leen las
mismas coordenadas al truncar cualquier prolongación posterior. La
conmutatividad es una igualdad de aplicaciones sobre estados producidos, no
la conservación tautológica de una salida incluida en la entrada.
\end{proof}

\section{Órbitas de paso tres y transformada de Hadamard}

El orden cíclico del sello se distribuye en las tres órbitas
\begin{align}
 K^{(1)}&=(K_1,K_4,K_7,K_{10}),
 \label{eq:orbita-k-1}\\
 K^{(2)}&=(K_2,K_5,K_8,K_{11}),
 \label{eq:orbita-k-2}\\
 K^{(3)}&=(K_3,K_6,K_9,K_{12}).
 \label{eq:orbita-k-3}
\end{align}
Sea
\begin{equation}
 H_4=
 \begin{pmatrix}
 1&1&1&1\\
 1&1&-1&-1\\
 1&-1&1&-1\\
 1&-1&-1&1
 \end{pmatrix}.
 \label{eq:matriz-hadamard-cuatro-autonoma}
\end{equation}
Si \(P_{\rm orb}\) reordena las doce coordenadas de acuerdo con
\eqref{eq:orbita-k-1}--\eqref{eq:orbita-k-3}, la transformación global es
\begin{equation}
 \mathcal H_{12}
 =P_{\rm orb}^{-1}
 (H_4\oplus H_4\oplus H_4)P_{\rm orb}.
 \label{eq:hadamard-global-doce}
\end{equation}

Los tres bloques de \eqref{eq:vector-u-firmado-autonomo}, leídos por
órbitas, son
\begin{align*}
 U^{(1)}&=(2378,-452,-668,-322),\\
 U^{(2)}&=(1406,998,-204,-28),\\
 U^{(3)}&=(2479,-551,-371,-997).
\end{align*}

\begin{lemma}[Cuarto de vuelta algebraico de Hadamard]
\label{lem:cuarto-inversa-hadamard}
La matriz \(H_4\) satisface
\begin{equation}
 H_4^2=4I_4,
 \qquad
 H_4^{-1}=\frac14H_4,
 \qquad
 \det H_4=-16.
 \label{eq:identidades-hadamard-cuatro}
\end{equation}
\end{lemma}

\begin{proof}
Las filas de \(H_4\) son ortogonales y cada una posee norma cuadrada cuatro.
Como la matriz es simétrica, \(H_4^{\mathsf T}H_4=H_4^2=4I_4\). La
fórmula de la inversa y el determinante se deducen inmediatamente.
\end{proof}

\begin{theorem}[Reconstrucción entera única del sello]
\label{thm:reconstruccion-entera-k-hadamard}
La inversión
\begin{equation}
 K^{(r)}=\frac14H_4U^{(r)}
 \qquad(r=1,2,3)
 \label{eq:inversion-bloques-hadamard}
\end{equation}
produce
\begin{align*}
 K^{(1)}&=(234,729,621,794),\\
 K^{(2)}&=(543,659,058,146),\\
 K^{(3)}&=(140,824,914,601).
\end{align*}
Al deshacer el orden orbital se obtiene
\begin{equation}
 \boxed{
 K=(234,543,140,729,659,824,621,058,914,794,146,601).}
 \label{eq:sello-k-doce-autonomo}
\end{equation}
Es la única preimagen racional de \(U\), y pertenece a
\(\mathbb Z^{12}\).
\end{theorem}

\begin{proof}
Por \eqref{eq:identidades-hadamard-cuatro}, la preimagen racional es única.
Las multiplicaciones explícitas son
\begin{align*}
 H_4U^{(1)}&=(936,2916,2484,3176),\\
 H_4U^{(2)}&=(2172,2636,232,584),\\
 H_4U^{(3)}&=(560,3296,3656,2404).
\end{align*}
Todas las coordenadas son divisibles por cuatro. Dividir y reintercalar
produce \eqref{eq:sello-k-doce-autonomo}.
\end{proof}

El factor \(1/4\) no se escoge para ajustar una salida: es la inversa
forzada por \(H_4^2=4I_4\). Tampoco es la raíz cuadrada de \(-1\). La
estructura compleja interna \(A_W^2=-I\) aparecerá después de construir la
matriz de Paley; son dos hechos de tipos diferentes.

\section{Estructura integral y forma normal de Smith}

\begin{theorem}[Imagen integral de \(H_4\)]
\label{thm:smith-hadamard-autonomo}
La forma normal de Smith es
\begin{equation}
 \operatorname{SNF}(H_4)=\operatorname{diag}(1,2,2,4).
 \label{eq:snf-hadamard-autonoma}
\end{equation}
En consecuencia,
\begin{equation}
 \operatorname{coker}\mathcal H_{12}
 \cong(\mathbb Z/2\mathbb Z)^6
 \oplus(\mathbb Z/4\mathbb Z)^3,
 \label{eq:cociente-hadamard-doce}
\end{equation}
y la imagen de \(\mathcal H_{12}\) posee índice
\[
 16^3=4096
\]
en \(\mathbb Z^{12}\). Para \(u\in\mathbb Z^4\),
\begin{equation}
 u\in H_4\mathbb Z^4
 \quad\Longleftrightarrow\quad
 H_4u\equiv0\pmod4.
 \label{eq:criterio-imagen-hadamard}
\end{equation}
\end{theorem}

\begin{proof}
El máximo común divisor de las entradas es uno. Los menores de orden dos son
pares y alguno tiene valor absoluto dos; los de orden tres son múltiplos de
cuatro y alguno tiene valor absoluto cuatro; el determinante tiene módulo
dieciséis. Los divisores determinantes son \(1,2,4,16\), de donde se obtienen
los factores invariantes \(1,2,2,4\). La suma directa triple repite los
factores y multiplica los índices. Finalmente,
\(H_4^{-1}=H_4/4\), por lo que la preimagen es entera exactamente bajo la
congruencia \eqref{eq:criterio-imagen-hadamard}.
\end{proof}

\section{Certificación en el segundo sistema de coordenadas: diferencias de pasos tres y cuatro}

Con índices cíclicos módulo doce, definimos
\begin{equation}
 (D_3K)_m=K_m-K_{m+3},
 \qquad
 (D_4K)_m=K_m-K_{m+4},
 \qquad
 Q(K)=\sum_{m=1}^{12}K_m.
 \label{eq:extractor-diferencias-k}
\end{equation}
Aplicados al sello ya reconstruido por Hadamard, estos operadores deben
recuperar exactamente las coordenadas de canal que el registro TPK produjo
antes de \(U\). Para \eqref{eq:sello-k-doce-autonomo},
\begin{align*}
 D_3K={}&(-495,-116,-684,108,601,-90,\\
         &\qquad-173,-88,313,560,-397,461),\\
 D_4K={}&(-425,-281,-481,671,-255,30,\\
         &\qquad475,-543,680,251,6,-128),\\
 Q(K)={}&6263.
\end{align*}
Por comparación con \eqref{eq:canales-ledger-previos-u}, queda probada la
igualdad de salidas
\begin{equation}
 \boxed{
 D_3K=b^{90},\qquad D_4K=b^{120},\qquad Q(K)=Q_{\rm TPK}.}
 \label{eq:certificacion-canales-ledger-desde-k}
\end{equation}
El sistema diferencial certifica y reconstruye el registro previo; no
selecciona \(K\), \(U\) ni ninguna constante.

\begin{theorem}[Completitud del extractor transversal]
\label{thm:completitud-extractor-transversal-k}
El operador
\[
 \mathcal E_{\rm dod}=(D_3,D_4,Q):
 \mathbb Z^{12}\longrightarrow\mathbb Z^{25}
\]
tiene rango doce y forma normal de Smith
\begin{equation}
 \operatorname{SNF}(\mathcal E_{\rm dod})
 =\operatorname{diag}(\underbrace{1,\ldots,1}_{11},12).
 \label{eq:snf-extractor-transversal-k}
\end{equation}
En particular, \((D_3K,D_4K,Q(K))\) determina un único \(K\).
\end{theorem}

\begin{proof}
Si \(D_3v=D_4v=0\), entonces \(v\) es constante a lo largo de los pasos
tres y cuatro. Como \(\gcd(3,4)=1\), esos pasos generan todo
\(\mathbb Z/12\mathbb Z\), de modo que
\[
 \ker D_3\cap\ker D_4=\langle\mathbf1\rangle.
\]
La carga total no se anula en esa recta, pues \(Q(\mathbf1)=12\). El rango
es doce. Las filas de \((D_3,D_4)\) son incidencias orientadas de un grafo
de Cayley conexo; un árbol generador aporta once factores invariantes
unitarios. Todo menor máximo debe contener la fila \(Q\), y las operaciones
unimodulares reducen su última contribución a doce. Existe un menor máximo de
módulo doce; por tanto el último factor es exactamente doce.
\end{proof}

La reconstrucción de \(U\) desde estas coordenadas no necesita volver a
multiplicar por \(H_4\). Para \(r=1,2,3\), sea
\begin{equation}
 B_r=\sum_{j=0}^{3}(D_4K)_{r+3j}.
 \label{eq:br-diferencias-k}
\end{equation}
Entonces
\[
 (B_1,B_2,B_3)=(972,-1073,101).
\]
Las sumas orbitales son
\begin{align}
 A_1&=\frac{Q+2B_1+B_2}{3}=2378,
 \label{eq:a1-reconstruccion-u}\\
 A_2&=A_1-B_1=1406,
 \label{eq:a2-reconstruccion-u}\\
 A_3&=A_2-B_2=2479.
 \label{eq:a3-reconstruccion-u}
\end{align}
Si \(d_{r,j}=(D_3K)_{r+3j}\), el bloque firmado de la órbita \(r\) es
\begin{equation}
 \bigl(A_r,
 d_{r,1}-d_{r,3},
 d_{r,0}+d_{r,2},
 d_{r,0}-d_{r,2}\bigr),
 \label{eq:bloque-firmado-desde-diferencias}
\end{equation}
que reproduce los tres bloques \(U^{(r)}\).

\section{Tétrada de umbral del sello}

La completación decimal local satisface
\[
 1000=729+271.
\]
El soporte de las coordenadas del sello que alcanzan el umbral \(729\) es
\begin{equation}
 \boxed{
 B_K:=\{m:K_m\geq729\}=\{4,6,9,10\}.}
 \label{eq:tetrada-umbral-k}
\end{equation}
Esta tétrada se obtiene después de construir \(K\). En esta unidad no se la
denomina todavía estrella de Witt: tal objeto sólo estará definido después
de construir el código ternario, sus hexadas y el sistema de Steiner.

\section{Independencia causal respecto de la salida física}

\begin{theorem}[Aislamiento de \(U\) y \(K\)]
\label{thm:aislamiento-u-k-alpha-codata}
La composición
\begin{equation}
 \Omega_{\rm seed}\longrightarrow x_{\rm term}
 \xrightarrow{\mathcal R_{12}}
 \mathcal R_{12}(x_{\rm term})
 \xrightarrow{\mathcal E_{108}^{90,120}}
 (b^{90},b^{120},Q_{\rm TPK})
 \xrightarrow{\Pi_H}U
 \xrightarrow{\mathcal H_{12}^{-1}}K
 \label{eq:composicion-upstream-u-k}
\end{equation}
no recibe como entrada \(\alpha\), \(\alpha^{-1}\), CODATA ni una cadena
decimal objetivo. El sello \(K\) queda determinado antes de definir la
recurrencia que producirá \(\alpha\).
\end{theorem}

\begin{proof}
La primera parte de la composición utiliza únicamente semillas, emisiones,
lectores ternarios, operadores \(L_0,L_1\), relaciones de supervivencia,
transporte de fase y las cartas del estado TPK. El registro
\(\mathcal R_{12}\) y el lector de sus canales producen
\((b^{90},b^{120},Q_{\rm TPK})\); las fórmulas
\eqref{eq:lector-armonico-a-desde-ledger} y
\eqref{eq:lector-armonico-bloque-desde-ledger} producen \(U\); y sólo
entonces actúa la transformación lineal
\(\mathcal H_{12}^{-1}=\mathcal H_{12}/4\) sobre su imagen integral.
Ninguna de las fórmulas contiene una variable física o una tabla
metrológica. La recurrencia en base mil se introduce sólo en la unidad
siguiente, después de haber fijado \(K\).
\end{proof}

\section{Criterios de refutación}

\begin{criterion}[Refutación del registro y del sello]
La construcción queda refutada si se verifica cualquiera de los hechos
siguientes:
\begin{enumerate}
 \item las ventanas \(E_m\) no particionan los ciento ocho pasos;
 \item la carta \(\mathcal L_{12}\) incumple su inversa o sus congruencias;
 \item el censo terminal no produce \(331\) ternas, dos matrices y una única
       carga axial;
 \item se intenta obtener \(U\) de \(B_{\rm term}\) después de olvidar las
       coordenadas firmadas;
 \item falla el cuadrado de naturalidad
       \eqref{eq:cuadrado-naturalidad-lector-firmado};
 \item \(H_4^2\neq4I_4\), alguna preimagen no es entera o el vector
       reconstruido difiere de \eqref{eq:sello-k-doce-autonomo};
 \item la forma normal de Smith difiere de
       \eqref{eq:snf-hadamard-autonoma};
 \item el extractor \((D_3,D_4,Q)\) pierde rango o no reproduce \(U\);
 \item \(B_K\) se introduce antes de construir \(K\), o se le atribuye una
       estructura de Witt antes de definir \(W_{12}\);
 \item una cifra de \(\alpha\) o una referencia metrológica interviene en
       cualquiera de las flechas de \eqref{eq:composicion-upstream-u-k}.
\end{enumerate}
\end{criterion}
